% !PRINTED 373
\weberpage{373}
Now one can, however, choose $b$ so that
\weq{17}{b^2 \equiv D \ (\mathrm{mod}\ 4a^n)}
holds, for an arbitrarily large $n$. For one can alter $b$ by an arbitrary multiple of $2a$. If we therefore assume (17) to be satisfied for an exponent $n$, and accordingly put
\weq{18}{b^2 = D + 4a^n c,}
then it follows that:
\[
(b + 2ha^n)^2 - D \equiv 4a^n(c+hb) \ (\mathrm{mod}\ 4a^{n+1}),
\]
and hence, if one determines $h$ from the congruence $c + hb \equiv 0 \ (\mathrm{mod}\ a)$, there results a $b$ satisfying the congruence (17) for $n+1$.

Since now $b$ is relatively prime to $a$, it follows by 2. that, if $n \geqq m$, then $a^m$ is the least natural number divisible by $\mathfrak{a}^m$, and from this there results the following rule of composition:

5. \gesperrt{If the form $(a, b, a^n c) = U$ satisfies the condition (18), and if $m \leqq n$, then}
\[
U^m = (a^m, b, a^{n-m}c).
\]
With this the classes of quadratic forms of the discriminant $\Delta$ are united into an \gesperrt{Abelian group} which is \gesperrt{isomorphic with the group of ideal classes}, and at the same time a way has been indicated by which one can actually carry out the composition in these groups by means of suitably chosen representatives\footnotemark[1]).

\begin{center}
\S\,102. \gesperrt{Composition of the Orders.}
\end{center}

An order $O = [Q]$ in the field $\Omega$ is completely determined by the conductor $Q$. Let $O_1$, $O_2$ be two orders with the conductors $Q_1$, $Q_2$, and let the greatest common divisor of $Q_1$ and $Q_2$ be $Q$. If then $M$ is the least common multiplum of $Q_1$ and $Q_2$, then
\weq{1}{Q_1 Q_2 = QM.}

\begingroup
\renewcommand{\thefootnote}{1)}
\footnotetext[1]{The \gesperrt{unit} in the group of classes is formed by the \gesperrt{principal class}. Concerning the composition of quadratic forms the following are to be mentioned: \gesperrt{Gauss}, \glqq Disquisitiones arithmeticae\grqq{}, Art. 235\,ff. \gesperrt{Dirichlet}, De formarum binariarum secundi gradus compositione (1851), Werke Bd. II, S. 105. \gesperrt{Dirichlet-Dedekind}, Vorlesungen über Zahlentheorie, 4. Aufl., \S\,145\,f. \gesperrt{Dedekind}, Crelles Journal, Bd. 129. \gesperrt{H. Weber}, Göttinger Nachrichten, 9. Februar 1907.}
\endgroup
